\begin{textsample}{POS Dim 3 – pb – Score 15.00 – pb/pb\_inf\_24.txt}  \label{ex:f3_pos_117}
CAMPO DE EXPERIÊNCIA : “ CORPO, \textbf{GESTOS} E MOVIMENTOS ” OBJETIVOS DE APRENDIZAGEM E DESENVOLVIMENTO Criar com o corpo formas diversificadas de expressão de sentimentos, sensações e emoções, tanto nas situações do cotidiano quanto em \textbf{brincadeiras}, dança, teatro e música. \textbf{Demonstrar} controle e adequação do uso do seu corpo em \textbf{brincadeiras} e jogos, \textbf{escuta} e reconto de histórias, atividades artísticas, entre outras possibilidades.
Criar movimentos, \textbf{gestos}, olhares e \textbf{mímicas} em \textbf{brincadeiras}, jogos e atividades artísticas, como dança, teatro e música.
Adotar \textbf{hábitos} de autocuidado relacionados à \textbf{higiene}, alimentação, \textbf{conforto} e aparência.
Coordenar suas habilidades manuais no atendimento adequado a seus interesses e necessidades em situações diversas. ( BNCC, 2017 ) SUGESTÕES DE VIVÊNCIAS QUE FAVORECEM A CONSTRUÇÃO DE EXPERIÊNCIAS PELAS \textbf{CRIANÇAS} \textbf{PEQUENAS} NESTE CAMPO As práticas realizadas na Instituição de Educação \textbf{Infantil} devem possibilitar : - Exploração de situações diversas que favoreçam experiências de expressão corporal por meio de diferentes linguagens ( dança, teatro, música ) e de \textbf{brincadeiras} do universo cultural da \textbf{criança} e de outras culturas ; - Situações em que possa vivenciar jogos, \textbf{brincadeiras}, atividades literárias, bem como outras, que impliquem o uso do seu corpo e um autocontrole do mesmo ; - Espaços de criação por meio de diferentes linguagens artísticas e que incluam o corpo e o movimento ; - Vivências de \textbf{cuidado} com o corpo, sua \textbf{higiene}, \textbf{conforto} e bem estar, construindo uma imagem de si a partir de sua própria estética, em termos de aparência ; - \textbf{Manipulem}, criem e confeccionem \textbf{brinquedos} da cultura local e geral ; - Possam vivenciar práticas relacionadas à autonomia, como cuidar do seu próprio corpo.

% matched lemmas: brincadeira, brinquedo, conforto, criança, cuidado, demonstrar, escuta, gesto, higiene, hábito, infantil, manipular, mímico, pequeno
\end{textsample}
